\documentclass{article}
\usepackage[margin=1in]{geometry}

% ==== PACKAGES ====
\usepackage{amsmath, amssymb, amsthm}
\usepackage{mathtools}
\usepackage[longnamesfirst,sort&compress]{natbib}
\usepackage{graphicx}
\usepackage{hyperref}
\usepackage{cleveref}
\usepackage{booktabs}
\usepackage{enumitem}

% Citation punctuation style
\bibpunct{(}{)}{;}{a}{,}{,}

% ==== CUSTOM COMMANDS ====
\newcommand{\E}{\mathbb{E}}
\renewcommand{\P}{\mathbb{P}}
\newcommand{\R}{\mathbb{R}}
\newcommand{\diff}{\mathrm{d}}
\newcommand{\F}{\mathcal{F}}
\newcommand{\one}{\mathbf{1}}
\newcommand{\Rea}{\operatorname{Re}}

% ==== THEOREM ENVIRONMENTS ====
\theoremstyle{definition}
\newtheorem{definition}{Definition}[section]
\newtheorem{remark}{Remark}[section]
\newtheorem{proposition}{Proposition}[section]
\newtheorem{lemma}{Lemma}[section]
\newtheorem{assumption}{Assumption}[section]

\title{Sepp-Style Laplace--Resolvent Semi-Analytics for a Ratio Barrier Payoff under Kou Dynamics}
\author{}
\date{\today}

\begin{document}
\maketitle

\tableofcontents

% ============================================================
\section{Problem Setup and One-Dimensional Reduction}
% ============================================================

Let $(\Omega,\F,(\F_t)_{t\ge 0},\P)$ be a filtered probability space.
Consider two strictly positive asset prices $S_1,S_2$ with initial values $S_{1,0},S_{2,0}>0$.
We write
\[
S_{i,t}=S_{i,0}e^{X_{i,t}},\qquad X_{i,0}=0,\quad i=1,2.
\]

Fix $b\in(0,1)$ and portfolio weights $w_1,w_2>0$. The health ratio is defined by
\[
H_t := \frac{b\,w_1S_{1,t}}{w_2S_{2,t}},
\qquad
h_t:=\log H_t,
\qquad
\tau:=\inf\{t\ge 0:\; H_t<1\}.
\]
Since
\[
h_t
=
\log\!\Big(\frac{b\,w_1S_{1,0}}{w_2S_{2,0}}\Big)+(X_{1,t}-X_{2,t}),
\]
it is natural to define
\[
h_0:=\log\!\Big(\frac{b\,w_1S_{1,0}}{w_2S_{2,0}}\Big),
\qquad
Z_t:=X_{1,t}-X_{2,t},
\]
so that
\[
h_t=h_0+Z_t,
\qquad
\tau=\inf\{t\ge 0:\; Z_t<-h_0\}.
\]

The terminal payoff is
\[
\Pi_T:=(w_1S_{1,T}-w_2S_{2,T})^+\,\one_{\{\tau>T\}}.
\]
Using
\[
H_T=\frac{b\,w_1S_{1,T}}{w_2S_{2,T}},
\]
one obtains the factorization
\begin{equation}\label{eq:payoff-factor}
\boxed{
\Pi_T
=
\frac{w_2}{b}\,S_{2,0}e^{X_{2,T}}\,(e^{h_0+Z_T}-b)^+\,\one_{\{\tau>T\}}.
}
\end{equation}

We impose the portfolio normalization
\[
w_1S_{1,0}-w_2S_{2,0}=1.
\]
This leaves only one degree of freedom. Since the valuation problem enters through $h_0$, it is natural to use
$h_0$ as the scalar control parameter. Let
\[
x:=w_2S_{2,0}.
\]
Then
\[
w_1S_{1,0}=1+x,
\qquad
e^{h_0}=\frac{b(1+x)}{x}.
\]
Solving for $x$ gives
\[
x=\frac{b}{e^{h_0}-b}.
\]
Hence
\begin{equation}\label{eq:w-from-h0}
\boxed{
w_2(h_0)=\frac{b}{S_{2,0}(e^{h_0}-b)},
\qquad
w_1(h_0)=\frac{e^{h_0}}{S_{1,0}(e^{h_0}-b)}.
}
\end{equation}
Since $w_1,w_2>0$, the feasible set is
\[
h_0>\ln b.
\]
Thus the original two-weight optimization problem is equivalently reduced to a one-dimensional optimization over
\[
h_0\in(\ln b,\infty).
\]

% ============================================================
\section{Kou Model, Tilted Dynamics, and Notation}
% ============================================================

We assume under $\P$ that $(X_{1,t},X_{2,t})_{t\ge 0}$ is a bivariate jump--diffusion with correlated Brownian part
(volatilities $\sigma_1,\sigma_2$, correlation $\rho$) and independent Kou compound Poisson jumps in each coordinate,
independent of the Brownian part.

Its joint L\'evy--Khintchine exponent is defined by
\[
\E_{\P}\!\left[e^{i(uX_{1,t}+vX_{2,t})}\right]=\exp\!\big(t\,\Psi(u,v)\big),\qquad u,v\in\R,
\]
with
\begin{equation}\label{eq:Psi}
\boxed{
\begin{aligned}
\Psi(u,v)
&= i(\mu_1 u+\mu_2 v)
-\tfrac12\Big(\sigma_1^2u^2+\sigma_2^2v^2+2\rho\sigma_1\sigma_2uv\Big)\\
&\quad +\lambda_1\big(\varphi_{J_1}(u)-1\big)
+\lambda_2\big(\varphi_{J_2}(v)-1\big).
\end{aligned}}
\end{equation}

The Kou jump law is
\[
f_i(j)
=
p_i\eta_{i,+}e^{-\eta_{i,+}j}\one_{\{j>0\}}
+
(1-p_i)\eta_{i,-}e^{\eta_{i,-}j}\one_{\{j<0\}},
\qquad
\eta_{i,+},\eta_{i,-}>0,
\quad
p_i\in(0,1),
\]
so that
\begin{equation}\label{eq:Kou-char-mgf}
\boxed{
\varphi_{J_i}(u)
=
p_i\frac{\eta_{i,+}}{\eta_{i,+}-iu}
+
(1-p_i)\frac{\eta_{i,-}}{\eta_{i,-}+iu},
\qquad
M_i(s):=\E[e^{sJ_i}]
=
p_i\frac{\eta_{i,+}}{\eta_{i,+}-s}
+
(1-p_i)\frac{\eta_{i,-}}{\eta_{i,-}+s}.
}
\end{equation}

For each integer $k\ge 1$, define the tilted measure $\P^{(2,k)}$ by
\[
\left.\frac{\diff\P^{(2,k)}}{\diff\P}\right|_{\F_t}
=
\frac{e^{kX_{2,t}}}{\E_{\P}[e^{kX_{2,t}}]}
=
\exp\!\left(kX_{2,t}-t\,\Psi(0,-ik)\right),
\]
assuming $\eta_{2,+}>k$. Under $\P^{(2,k)}$, the one-dimensional process $Z$ is again L\'evy with exponent
\begin{equation}\label{eq:psiZk}
\boxed{
\psi_Z^{(k)}(s)
=
\mu_Z^{(k)}s+\tfrac12\sigma_Z^2 s^2
+\lambda_1\big(M_1(s)-1\big)
+\lambda_2\big(M_2(k-s)-M_2(k)\big),
}
\end{equation}
where
\[
\sigma_Z^2:=\sigma_1^2+\sigma_2^2-2\rho\sigma_1\sigma_2 \ge 0,
\qquad
\mu_Z^{(k)}:=(\mu_1-\mu_2)-k(\sigma_2^2-\rho\sigma_1\sigma_2).
\]
Let $\mathcal L^{(k)}$ denote the generator of $Z$ under $\P^{(2,k)}$. Then
\begin{equation}\label{eq:generator-eigen}
\boxed{
\mathcal L^{(k)}e^{sz}=\psi_Z^{(k)}(s)e^{sz}.
}
\end{equation}

For the Kou structure under the $k$-tilted measure, we introduce the $k$-dependent phase rates
\[
r_{1,+}^{(k)}:=\eta_{1,+},
\qquad
r_{2,+}^{(k)}:=\eta_{2,-}+k,
\]
\[
r_{1,-}^{(k)}:=\eta_{1,-},
\qquad
r_{2,-}^{(k)}:=\eta_{2,+}-k.
\]
The associated phase operators are
\[
K_{j,+}^{(k)}[f](z)
:=
\int_0^\infty f(z+y)\,r_{j,+}^{(k)}e^{-r_{j,+}^{(k)}y}\,\diff y,
\qquad j=1,2,
\]
and
\[
K_{j,-}^{(k)}[f](z)
:=
\int_{-\infty}^0 f(z+y)\,r_{j,-}^{(k)}e^{r_{j,-}^{(k)}y}\,\diff y,
\qquad j=1,2.
\]
For exponentials,
\begin{equation}\label{eq:phase-on-exp}
\boxed{
K_{j,+}^{(k)}[e^{\gamma\cdot}](z)
=
\frac{r_{j,+}^{(k)}}{r_{j,+}^{(k)}-\gamma}e^{\gamma z},
\qquad
K_{j,-}^{(k)}[e^{\gamma\cdot}](z)
=
\frac{r_{j,-}^{(k)}}{r_{j,-}^{(k)}+\gamma}e^{\gamma z}.
}
\end{equation}

The role of these phase rates is structural: they convert the nonlocal jump terms into a finite set of phase variables.
This is what makes the matching conditions at the payoff kink and the absorbing conditions at the barrier reducible to
finite-dimensional linear systems.

% ============================================================
\section{Moment Reduction}
% ============================================================

Define
\[
g_k(z;h_0):=\big((e^{h_0+z}-b)^+\big)^k.
\]
From \eqref{eq:payoff-factor},
\[
\Pi_T^k
=
\Big(\frac{w_2}{b}\Big)^k S_{2,0}^k e^{kX_{2,T}}\,g_k(Z_T;h_0)\,\one_{\{\tau>T\}}.
\]
Let
\[
A_k:=\E_{\P}[e^{kX_{2,T}}]=\exp\!\big(T\Psi(0,-ik)\big).
\]
Then
\begin{equation}\label{eq:moment-reduction}
\boxed{
\E_{\P}[\Pi_T^k]
=
\Big(\frac{w_2}{b}\Big)^k S_{2,0}^k\,A_k\;
\E_{\P^{(2,k)}}\!\Big[g_k(Z_T;h_0)\,\one_{\{\tau>T\}}\Big].
}
\end{equation}

\begin{proposition}[First two moments]\label{prop:first-two-moments}
Assume $k=1,2$ satisfy $\eta_{2,+}>k$. Define
\[
m_k(T;h_0):=\E_{\P^{(2,k)}}\!\Big[g_k(Z_T;h_0)\,\one_{\{\tau>T\}}\Big].
\]
Then
\[
\E_{\P}[\Pi_T]=\Big(\frac{w_2}{b}\Big) S_{2,0}\,e^{T\Psi(0,-i)}\,m_1(T;h_0),
\]
\[
\E_{\P}[\Pi_T^2]=\Big(\frac{w_2}{b}\Big)^2 S_{2,0}^2\,e^{T\Psi(0,-2i)}\,m_2(T;h_0),
\]
and consequently
\[
\mathrm{Var}_{\P}(\Pi_T)=\E_{\P}[\Pi_T^2]-\big(\E_{\P}[\Pi_T]\big)^2.
\]
\end{proposition}

The remainder of the analysis is therefore devoted to computing the killed expectations $m_1(T;h_0)$ and $m_2(T;h_0)$.

% ============================================================
\section{Three-Step Resolvent Construction}
% ============================================================

The valuation problem is most transparent when built progressively. We begin with the atomic response to a single
exponential source, then restore the payoff kink on the full line, and finally impose the absorbing barrier by a
homogeneous correction.

\paragraph{Step 1: atomic full-line resolvent.}
The first object is the Laplace transform of exponential moments,
\[
\widehat U_{\mathrm{atom}}^{(k)}(q,z;s)
=
\int_0^\infty e^{-qt}\,\E_z^{(2,k)}[e^{sZ_t}]\,\diff t.
\]
This is not yet the payoff value itself, but it is the basic building block because the payoff-active branch is a
finite linear combination of exponentials.

The problem is
\[
(q-\mathcal L^{(k)})u(z)=e^{sz},\qquad z\in\R.
\]
The domain is the full line $\R$; there is no interface and no boundary.

\begin{proposition}[Step 1: atomic resolvent]\label{prop:Uatom}
For any $q>0$ and any $s$ in the moment domain such that $q\neq \psi_Z^{(k)}(s)$,
\[
\boxed{
\widehat U_{\mathrm{atom}}^{(k)}(q,z;s)
=
\frac{e^{sz}}{q-\psi_Z^{(k)}(s)}.
}
\]
Equivalently,
\[
(q-\mathcal L^{(k)})\widehat U_{\mathrm{atom}}^{(k)}(q,z;s)=e^{sz}.
\]
\end{proposition}

\paragraph{Step 2: full-line resolvent with payoff kink.}
We now restore the actual payoff but still ignore liquidation. The quantity of interest is
\[
\widehat U_{\mathrm{kink}}^{(k)}(q,z;h_0)
=
\int_0^\infty e^{-qt}
\E_z^{(2,k)}\!\left[\big(e^{h_0+Z_t}-b\big)_+^k\right]\diff t.
\]
This is the no-barrier payoff value.

The problem is
\[
(q-\mathcal L^{(k)})\widehat U_{\mathrm{kink}}^{(k)}(q,z;h_0)
=
g_k(z;h_0),
\qquad z\in\R,
\]
with
\[
g_k(z;h_0)=\big(e^{h_0+z}-b\big)_+^k.
\]
Define
\[
z_*(h_0):=\ln b-h_0.
\]
Then
\[
g_k(z;h_0)=
\begin{cases}
0, & z<z_*(h_0),\\[4pt]
(e^{h_0+z}-b)^k, & z>z_*(h_0).
\end{cases}
\]
For $z>z_*(h_0)$,
\[
(e^{h_0+z}-b)^k
=
\sum_{j=0}^k c_{k,j}^{(z)}(h_0)e^{jz},
\qquad
c_{k,j}^{(z)}(h_0):=\binom{k}{j}(-1)^{k-j}b^{\,k-j}e^{jh_0}.
\]

The domain is the full line $\R$. There is a single interface at $z=z_*(h_0)$. There is no barrier boundary in this
step; instead, one imposes asymptotic boundedness as $z\to\pm\infty$.

\begin{assumption}[Generic six-root regime]\label{ass:six-root}
Fix $k\in\{1,2\}$ and $q>0$. Assume:
\begin{enumerate}[label=(\roman*),leftmargin=2em]
\item $\sigma_Z>0$;
\item the characteristic equation
\[
\psi_Z^{(k)}(\gamma)=q
\]
has six simple roots $\gamma_1(q),\dots,\gamma_6(q)$, ordered so that
\[
\Rea(\gamma_{1,2,3})>0,\qquad \Rea(\gamma_{4,5,6})<0;
\]
\item $q\neq \psi_Z^{(k)}(j)$ for all $j=0,\dots,k$.
\end{enumerate}
\end{assumption}

\begin{proposition}[Step 2: full-line kinked resolvent]\label{prop:Ukink}
Under \Cref{ass:six-root}, the full-line resolvent admits the decomposition
\[
\widehat U_{\mathrm{kink}}^{(k)}(q,z;h_0)
=
\widehat U_{\mathrm{pay}}^{(k)}(q,z;h_0)
+
\widehat U_{\mathrm{kink\text{-}corr}}^{(k)}(q,z;h_0),
\]
where the payoff-driven particular term is
\[
\boxed{
\widehat U_{\mathrm{pay}}^{(k)}(q,z;h_0)
=
\one_{\{z>z_*(h_0)\}}
\sum_{j=0}^{k}
\frac{c_{k,j}^{(z)}(h_0)}{q-\psi_Z^{(k)}(j)}e^{jz}.
}
\]
Moreover,
\[
\boxed{
\widehat U_{\mathrm{kink}}^{(k)}(q,z;h_0)=
\begin{cases}
\displaystyle \sum_{m=1}^3 A_m^{(k)}(q;h_0)e^{\gamma_m(q)z}, & z<z_*(h_0),\\[10pt]
\displaystyle \sum_{m=4}^6 A_m^{(k)}(q;h_0)e^{\gamma_m(q)z}
+\sum_{j=0}^{k}\frac{c_{k,j}^{(z)}(h_0)}{q-\psi_Z^{(k)}(j)}e^{jz}, & z>z_*(h_0).
\end{cases}
}
\]

The coefficients $A_1^{(k)},\dots,A_6^{(k)}$ are determined by the following interface conditions at $z=z_*(h_0)$:
\begin{enumerate}[label=(\roman*),leftmargin=2em]
\item continuity of the value;
\item continuity of the first derivative;
\item continuity of the downward phase variables
\[
K_{j,-}^{(k)}[\widehat U_{\mathrm{kink}}^{(k)}](z_*(h_0)-)
=
K_{j,-}^{(k)}[\widehat U_{\mathrm{kink}}^{(k)}](z_*(h_0)+),
\qquad j=1,2;
\]
\item continuity of the upward phase variables
\[
K_{j,+}^{(k)}[\widehat U_{\mathrm{kink}}^{(k)}](z_*(h_0)-)
=
K_{j,+}^{(k)}[\widehat U_{\mathrm{kink}}^{(k)}](z_*(h_0)+),
\qquad j=1,2.
\]
\end{enumerate}

Equivalently, defining
\[
\mathbf A^{(k)}(q;h_0):=
\begin{pmatrix}
A_1^{(k)}(q;h_0)\\
A_2^{(k)}(q;h_0)\\
A_3^{(k)}(q;h_0)\\
A_4^{(k)}(q;h_0)\\
A_5^{(k)}(q;h_0)\\
A_6^{(k)}(q;h_0)
\end{pmatrix},
\]
\[
P_0^{(k)}(q;h_0):=
\sum_{m=0}^{k}\frac{c_{k,m}^{(z)}(h_0)}{q-\psi_Z^{(k)}(m)}e^{mz_*(h_0)},
\]
\[
P_1^{(k)}(q;h_0):=
\sum_{m=0}^{k}m\,\frac{c_{k,m}^{(z)}(h_0)}{q-\psi_Z^{(k)}(m)}e^{mz_*(h_0)},
\]
and, for $j=1,2$,
\[
P_{j,-}^{(k)}(q;h_0):=
\sum_{m=0}^{k}\frac{c_{k,m}^{(z)}(h_0)}{q-\psi_Z^{(k)}(m)}
\frac{r_{j,-}^{(k)}}{r_{j,-}^{(k)}+m}e^{mz_*(h_0)},
\]
\[
P_{j,+}^{(k)}(q;h_0):=
\sum_{m=0}^{k}\frac{c_{k,m}^{(z)}(h_0)}{q-\psi_Z^{(k)}(m)}
\frac{r_{j,+}^{(k)}}{r_{j,+}^{(k)}-m}e^{mz_*(h_0)},
\]
and
\[
\mathbf p^{(k)}(q;h_0):=
\begin{pmatrix}
P_0^{(k)}(q;h_0)\\
P_1^{(k)}(q;h_0)\\
P_{1,-}^{(k)}(q;h_0)\\
P_{2,-}^{(k)}(q;h_0)\\
P_{1,+}^{(k)}(q;h_0)\\
P_{2,+}^{(k)}(q;h_0)
\end{pmatrix},
\]
the coefficients solve
\[
\mathbf M_{\mathrm{kink}}^{(k)}(q;h_0)\,\mathbf A^{(k)}(q;h_0)=\mathbf p^{(k)}(q;h_0),
\]
where
\[
\mathbf M_{\mathrm{kink}}^{(k)}(q;h_0)=
\begin{pmatrix}
e^{\gamma_1 z_*} & e^{\gamma_2 z_*} & e^{\gamma_3 z_*} & -e^{\gamma_4 z_*} & -e^{\gamma_5 z_*} & -e^{\gamma_6 z_*}\\
\gamma_1 e^{\gamma_1 z_*} & \gamma_2 e^{\gamma_2 z_*} & \gamma_3 e^{\gamma_3 z_*}
& -\gamma_4 e^{\gamma_4 z_*} & -\gamma_5 e^{\gamma_5 z_*} & -\gamma_6 e^{\gamma_6 z_*}\\
\frac{r_{1,-}^{(k)}}{r_{1,-}^{(k)}+\gamma_1}e^{\gamma_1 z_*} &
\frac{r_{1,-}^{(k)}}{r_{1,-}^{(k)}+\gamma_2}e^{\gamma_2 z_*} &
\frac{r_{1,-}^{(k)}}{r_{1,-}^{(k)}+\gamma_3}e^{\gamma_3 z_*} &
-\frac{r_{1,-}^{(k)}}{r_{1,-}^{(k)}+\gamma_4}e^{\gamma_4 z_*} &
-\frac{r_{1,-}^{(k)}}{r_{1,-}^{(k)}+\gamma_5}e^{\gamma_5 z_*} &
-\frac{r_{1,-}^{(k)}}{r_{1,-}^{(k)}+\gamma_6}e^{\gamma_6 z_*}\\
\frac{r_{2,-}^{(k)}}{r_{2,-}^{(k)}+\gamma_1}e^{\gamma_1 z_*} &
\frac{r_{2,-}^{(k)}}{r_{2,-}^{(k)}+\gamma_2}e^{\gamma_2 z_*} &
\frac{r_{2,-}^{(k)}}{r_{2,-}^{(k)}+\gamma_3}e^{\gamma_3 z_*} &
-\frac{r_{2,-}^{(k)}}{r_{2,-}^{(k)}+\gamma_4}e^{\gamma_4 z_*} &
-\frac{r_{2,-}^{(k)}}{r_{2,-}^{(k)}+\gamma_5}e^{\gamma_5 z_*} &
-\frac{r_{2,-}^{(k)}}{r_{2,-}^{(k)}+\gamma_6}e^{\gamma_6 z_*}\\
\frac{r_{1,+}^{(k)}}{r_{1,+}^{(k)}-\gamma_1}e^{\gamma_1 z_*} &
\frac{r_{1,+}^{(k)}}{r_{1,+}^{(k)}-\gamma_2}e^{\gamma_2 z_*} &
\frac{r_{1,+}^{(k)}}{r_{1,+}^{(k)}-\gamma_3}e^{\gamma_3 z_*} &
-\frac{r_{1,+}^{(k)}}{r_{1,+}^{(k)}-\gamma_4}e^{\gamma_4 z_*} &
-\frac{r_{1,+}^{(k)}}{r_{1,+}^{(k)}-\gamma_5}e^{\gamma_5 z_*} &
-\frac{r_{1,+}^{(k)}}{r_{1,+}^{(k)}-\gamma_6}e^{\gamma_6 z_*}\\
\frac{r_{2,+}^{(k)}}{r_{2,+}^{(k)}-\gamma_1}e^{\gamma_1 z_*} &
\frac{r_{2,+}^{(k)}}{r_{2,+}^{(k)}-\gamma_2}e^{\gamma_2 z_*} &
\frac{r_{2,+}^{(k)}}{r_{2,+}^{(k)}-\gamma_3}e^{\gamma_3 z_*} &
-\frac{r_{2,+}^{(k)}}{r_{2,+}^{(k)}-\gamma_4}e^{\gamma_4 z_*} &
-\frac{r_{2,+}^{(k)}}{r_{2,+}^{(k)}-\gamma_5}e^{\gamma_5 z_*} &
-\frac{r_{2,+}^{(k)}}{r_{2,+}^{(k)}-\gamma_6}e^{\gamma_6 z_*}
\end{pmatrix},
\qquad z_*=z_*(h_0).
\]
\end{proposition}

\paragraph{Step 3: killed problem with absorbing barrier.}
We finally restore liquidation. The resulting object is
\[
\widehat U_{\mathrm{bar}}^{(k)}(q,z;h_0)
=
\int_0^\infty e^{-qt}
\E_z^{(2,k)}\!\left[\big(e^{h_0+Z_t}-b\big)_+^k\one_{\{\tau(h_0)>t\}}\right]\diff t.
\]
This is the final barrier-adjusted payoff value.

The problem is
\[
(q-\mathcal L^{(k)})\widehat U_{\mathrm{bar}}^{(k)}(q,z;h_0)
=
g_k(z;h_0),
\qquad z>-h_0,
\]
with
\[
\widehat U_{\mathrm{bar}}^{(k)}(q,z;h_0)=0,
\qquad z\le -h_0.
\]

The domain is the half-line $(-h_0,\infty)$. The boundary is absorbing at $z=-h_0$. There is no new unknown interface;
the kink was already resolved in Step 2. However, because $z_*(h_0)<-h_0$, the downward convolution terms evaluated at
the barrier cross the kink.

\begin{proposition}[Step 3: barrier correction]\label{prop:Ubar}
Under \Cref{ass:six-root}, the killed resolvent may be written as
\[
\widehat U_{\mathrm{bar}}^{(k)}(q,z;h_0)
=
\widehat U_{\mathrm{kink}}^{(k)}(q,z;h_0)
+
\widehat U_{\mathrm{corr}}^{(k)}(q,z;h_0),
\qquad z>-h_0,
\]
where the correction is homogeneous:
\[
(q-\mathcal L^{(k)})\widehat U_{\mathrm{corr}}^{(k)}(q,z;h_0)=0,
\qquad z>-h_0.
\]
Moreover,
\[
\boxed{
\widehat U_{\mathrm{corr}}^{(k)}(q,z;h_0)
=
B_1^{(k)}(q;h_0)e^{\gamma_4(q)(z+h_0)}
+
B_2^{(k)}(q;h_0)e^{\gamma_5(q)(z+h_0)}
+
B_3^{(k)}(q;h_0)e^{\gamma_6(q)(z+h_0)}.
}
\]
The coefficients are determined by the barrier conditions
\[
\widehat U_{\mathrm{bar}}^{(k)}(q,-h_0^+;h_0)=0,
\]
\[
K_{1,-}^{(k)}\!\left[\widehat U_{\mathrm{bar}}^{(k)}(q,\cdot;h_0)\right](-h_0^+)=0,
\qquad
K_{2,-}^{(k)}\!\left[\widehat U_{\mathrm{bar}}^{(k)}(q,\cdot;h_0)\right](-h_0^+)=0.
\]
Equivalently,
\[
\mathbf M_{\mathrm{bar}}^{(k)}(q)\,\mathbf B^{(k)}(q;h_0)
=
-\mathbf b^{(k)}(q;h_0),
\]
where
\[
\mathbf B^{(k)}(q;h_0):=
\begin{pmatrix}
B_1^{(k)}(q;h_0)\\
B_2^{(k)}(q;h_0)\\
B_3^{(k)}(q;h_0)
\end{pmatrix},
\qquad
\mathbf b^{(k)}(q;h_0):=
\begin{pmatrix}
u_0^{(k)}(q;h_0)\\
u_{1,-}^{(k)}(q;h_0)\\
u_{2,-}^{(k)}(q;h_0)
\end{pmatrix},
\]
with
\[
u_0^{(k)}(q;h_0)=\widehat U_{\mathrm{kink}}^{(k)}(q,-h_0;h_0),
\]
\[
u_{j,-}^{(k)}(q;h_0)
=
K_{j,-}^{(k)}\!\left[\widehat U_{\mathrm{kink}}^{(k)}(q,\cdot;h_0)\right](-h_0),
\qquad j=1,2.
\]
The barrier matrix is
\[
\boxed{
\mathbf M_{\mathrm{bar}}^{(k)}(q)=
\begin{pmatrix}
1 & 1 & 1\\[6pt]
\dfrac{r_{1,-}^{(k)}}{r_{1,-}^{(k)}+\gamma_4(q)} &
\dfrac{r_{1,-}^{(k)}}{r_{1,-}^{(k)}+\gamma_5(q)} &
\dfrac{r_{1,-}^{(k)}}{r_{1,-}^{(k)}+\gamma_6(q)}\\[12pt]
\dfrac{r_{2,-}^{(k)}}{r_{2,-}^{(k)}+\gamma_4(q)} &
\dfrac{r_{2,-}^{(k)}}{r_{2,-}^{(k)}+\gamma_5(q)} &
\dfrac{r_{2,-}^{(k)}}{r_{2,-}^{(k)}+\gamma_6(q)}
\end{pmatrix}.
}
\]
Assuming invertibility,
\[
\mathbf B^{(k)}(q;h_0)
=
-\big(\mathbf M_{\mathrm{bar}}^{(k)}(q)\big)^{-1}\mathbf b^{(k)}(q;h_0).
\]
\end{proposition}

\begin{remark}[Explicit boundary inputs for $b\in(0,1)$]
Since the barrier point lies on the right branch of $\widehat U_{\mathrm{kink}}^{(k)}$ but the downward convolutions
cross the kink,
\[
u_0^{(k)}(q;h_0)
=
\sum_{m=4}^6 A_m^{(k)}(q;h_0)e^{-\gamma_m(q)h_0}
+
\sum_{j=0}^{k}\frac{c_{k,j}^{(z)}(h_0)}{q-\psi_Z^{(k)}(j)}e^{-jh_0},
\]
and for $j=1,2$,
\[
u_{j,-}^{(k)}(q;h_0)
=
\sum_{m=1}^3
A_m^{(k)}(q;h_0)e^{-\gamma_m(q)h_0}
\frac{r_{j,-}^{(k)}}{r_{j,-}^{(k)}+\gamma_m(q)}
\,b^{\,r_{j,-}^{(k)}+\gamma_m(q)}
\]
\[
\qquad
+\sum_{m=4}^6
A_m^{(k)}(q;h_0)e^{-\gamma_m(q)h_0}
\frac{r_{j,-}^{(k)}}{r_{j,-}^{(k)}+\gamma_m(q)}
\Big(1-b^{\,r_{j,-}^{(k)}+\gamma_m(q)}\Big)
\]
\[
\qquad
+\sum_{m=0}^{k}
\frac{c_{k,m}^{(z)}(h_0)}{q-\psi_Z^{(k)}(m)}e^{-mh_0}
\frac{r_{j,-}^{(k)}}{r_{j,-}^{(k)}+m}
\Big(1-b^{\,r_{j,-}^{(k)}+m}\Big).
\]
\end{remark}

% ============================================================
\section{Probability of Liquidation}
% ============================================================

The same resolvent architecture yields the probability of liquidation. Define
\[
p_{\mathrm{surv}}(t,z;h_0):=\P_z(\tau(h_0)>t),
\qquad
p_{\mathrm{liq}}(t,z;h_0):=1-p_{\mathrm{surv}}(t,z;h_0).
\]
Let $\mathcal L^{(0)}$ denote the generator of $Z$ under the untitled measure, with exponent $\psi_Z^{(0)}$, and define
\[
\widehat U_{\mathrm{surv}}(q,z;h_0)
:=
\int_0^\infty e^{-qt}p_{\mathrm{surv}}(t,z;h_0)\,\diff t.
\]

The problem is
\[
(q-\mathcal L^{(0)})\widehat U_{\mathrm{surv}}(q,z;h_0)=1,
\qquad z>-h_0,
\]
with
\[
\widehat U_{\mathrm{surv}}(q,z;h_0)=0,\qquad z\le -h_0.
\]
The domain is the half-line $(-h_0,\infty)$, the boundary is absorbing at $z=-h_0$, and there is no interface because
the source is constant.

\begin{proposition}[Survival resolvent and liquidation probability]\label{prop:survival}
Let $\gamma_4^{(0)}(q),\gamma_5^{(0)}(q),\gamma_6^{(0)}(q)$ be the three roots of
\[
\psi_Z^{(0)}(\gamma)=q
\]
with negative real part. Then
\[
\boxed{
\widehat U_{\mathrm{surv}}(q,z;h_0)
=
\frac{1}{q}
+
C_1(q;h_0)e^{\gamma_4^{(0)}(q)(z+h_0)}
+
C_2(q;h_0)e^{\gamma_5^{(0)}(q)(z+h_0)}
+
C_3(q;h_0)e^{\gamma_6^{(0)}(q)(z+h_0)}.
}
\]
The coefficients are determined by
\[
\widehat U_{\mathrm{surv}}(q,-h_0^+;h_0)=0,
\qquad
K_{1,-}^{(0)}[\widehat U_{\mathrm{surv}}](-h_0^+)=0,
\qquad
K_{2,-}^{(0)}[\widehat U_{\mathrm{surv}}](-h_0^+)=0.
\]
Equivalently,
\[
\mathbf M_{\mathrm{surv}}(q)\,\mathbf C(q;h_0)
=
-
\begin{pmatrix}
1/q\\
1/q\\
1/q
\end{pmatrix},
\]
where
\[
\mathbf C(q;h_0):=
\begin{pmatrix}
C_1(q;h_0)\\
C_2(q;h_0)\\
C_3(q;h_0)
\end{pmatrix},
\]
and
\[
\mathbf M_{\mathrm{surv}}(q)=
\begin{pmatrix}
1 & 1 & 1\\[6pt]
\dfrac{r_{1,-}^{(0)}}{r_{1,-}^{(0)}+\gamma_4^{(0)}(q)} &
\dfrac{r_{1,-}^{(0)}}{r_{1,-}^{(0)}+\gamma_5^{(0)}(q)} &
\dfrac{r_{1,-}^{(0)}}{r_{1,-}^{(0)}+\gamma_6^{(0)}(q)}\\[12pt]
\dfrac{r_{2,-}^{(0)}}{r_{2,-}^{(0)}+\gamma_4^{(0)}(q)} &
\dfrac{r_{2,-}^{(0)}}{r_{2,-}^{(0)}+\gamma_5^{(0)}(q)} &
\dfrac{r_{2,-}^{(0)}}{r_{2,-}^{(0)}+\gamma_6^{(0)}(q)}
\end{pmatrix}.
\]
At the initial state $Z_0=0$,
\[
\widehat U_{\mathrm{surv}}(q,0;h_0)
=
\frac{1}{q}
+
C_1(q;h_0)e^{\gamma_4^{(0)}(q)h_0}
+
C_2(q;h_0)e^{\gamma_5^{(0)}(q)h_0}
+
C_3(q;h_0)e^{\gamma_6^{(0)}(q)h_0}.
\]
Hence
\[
p_{\mathrm{surv}}(T,0;h_0)
=
\mathcal L^{-1}_{q\to T}\!\left[\widehat U_{\mathrm{surv}}(q,0;h_0)\right],
\qquad
p_{\mathrm{liq}}(T,0;h_0)=1-p_{\mathrm{surv}}(T,0;h_0).
\]
\end{proposition}

% ============================================================
\section{Bromwich Inversion and Numerical Stability}
% ============================================================

After the previous sections, the relevant Laplace-domain objects are
\[
\widehat U_{\mathrm{bar}}^{(k)}(q,z;h_0)
\qquad\text{and}\qquad
\widehat U_{\mathrm{surv}}(q,z;h_0).
\]
They are inverted in time via the Bromwich formula:
\[
U^{(k)}(t,z;h_0)
=
\frac{1}{2\pi i}
\int_{\sigma-i\infty}^{\sigma+i\infty}
e^{qt}\widehat U_{\mathrm{bar}}^{(k)}(q,z;h_0)\,\diff q,
\]
\[
p_{\mathrm{surv}}(t,z;h_0)
=
\frac{1}{2\pi i}
\int_{\sigma-i\infty}^{\sigma+i\infty}
e^{qt}\widehat U_{\mathrm{surv}}(q,z;h_0)\,\diff q,
\]
where $\sigma>0$ is the Bromwich shift.

The explicit formulas above are derived in the generic regime where the characteristic equation
\[
\psi_Z^{(k)}(\gamma)=q
\]
has six simple roots and there is no resonance at the forcing exponents:
\[
q\neq \psi_Z^{(k)}(j),\qquad j=0,\dots,k.
\]
Repeated roots occur at values satisfying
\[
\psi_Z^{(k)}(\gamma_\star)=q,
\qquad
(\psi_Z^{(k)})'(\gamma_\star)=0.
\]
These exceptional values are isolated in the generic regime.

A practical Bromwich shift $\sigma$ should therefore be chosen sufficiently positive so that the inversion nodes
\[
q=\sigma+i\omega
\]
remain well inside the generic six-root regime and away from
\begin{itemize}[leftmargin=2em]
\item resonance values $q=\psi_Z^{(k)}(j)$,
\item repeated-root values $q=\psi_Z^{(k)}(\gamma_\star)$ with $(\psi_Z^{(k)})'(\gamma_\star)=0$.
\end{itemize}

For each numerical node $q$, one should monitor
\[
\min_{j=0,\dots,k}|q-\psi_Z^{(k)}(j)|,
\qquad
\min_{a\neq b}|\gamma_a(q)-\gamma_b(q)|,
\]
and the condition numbers
\[
\kappa\!\left(\mathbf M_{\mathrm{kink}}^{(k)}(q;h_0)\right),
\qquad
\kappa\!\left(\mathbf M_{\mathrm{bar}}^{(k)}(q)\right).
\]

If repeated roots occur, the homogeneous basis must be enlarged to
\[
e^{\gamma z},\quad z e^{\gamma z},\quad \dots,\quad z^{m-1}e^{\gamma z}
\]
for a root of multiplicity $m$. If resonance occurs at a forcing exponent $j$, so that
\[
q=\psi_Z^{(k)}(j),
\]
the term
\[
\frac{c_{k,j}^{(z)}(h_0)}{q-\psi_Z^{(k)}(j)}e^{jz}
\]
must be replaced by a generalized particular ansatz of the form
\[
z^\nu e^{jz},
\]
with the minimal power $\nu$ restoring linear independence from the homogeneous solution space.

% ============================================================
\section{Consequence for Optimization}
% ============================================================

Because of \eqref{eq:w-from-h0}, every performance criterion can be written as a scalar function of $h_0$.
In particular,
\[
\E[\Pi_T^k]
=
\Big(\frac{w_2(h_0)}{b}\Big)^k S_{2,0}^k e^{T\Psi(0,-ik)}\,m_k(T;h_0)
=
\frac{e^{T\Psi(0,-ik)}}{(e^{h_0}-b)^k}\,m_k(T;h_0),
\]
where
\[
m_k(T;h_0)=U^{(k)}(T,0;h_0).
\]
Likewise,
\[
p_{\mathrm{liq}}(T,0;h_0)=1-p_{\mathrm{surv}}(T,0;h_0).
\]

Thus the portfolio optimization reduces to a one-dimensional problem over
\[
h_0\in(\ln b,\infty),
\]
after which the optimal weights are recovered from \eqref{eq:w-from-h0}.

% ============================================================
\section{Bibliography}
% ============================================================
\bibliographystyle{abbrvnamed}
\bibliography{finance}

\end{document}